% Generated draft: requires mathematical and source review.
\documentclass[11pt]{article}
\usepackage[margin=25mm]{geometry}
\usepackage{fontspec}
\setmainfont{Latin Modern Roman}
\usepackage{amsmath,amssymb,mathtools,unicode-math}
\setmathfont{Latin Modern Math}
\usepackage{enumitem,needspace,xurl,hyperref}
\setlength{\parindent}{0pt}
\setlength{\parskip}{5pt}
\setlength{\emergencystretch}{3em}
\newcommand{\N}{\mathbb{N}}
\newcommand{\Z}{\mathbb{Z}}
\newcommand{\Q}{\mathbb{Q}}
\newcommand{\R}{\mathbb{R}}
\newcommand{\C}{\mathbb{C}}
\newcommand{\F}{\mathbb{F}}
\newcommand{\A}{\mathbb{A}}
\newcommand{\PP}{\mathbb{P}}
\newcommand{\E}{\mathbb{E}}
\newcommand{\eps}{\varepsilon}
\newcommand{\dd}{\,\mathrm{d}}
\begin{document}
\section*{UnsolvedMath ID: 700419}
\textbf{Source problem code:} catalogue entry
\subsection{Definitions and mathematical statement}
Let \(k\geq 1\) and \(m\geq 0\) be integers, let \(X=\{x_1,\ldots,x_k\}\), and let \(F(X)\) be the free group on \(X\). For relators \(r_1,\ldots,r_m\in F(X)\), put
\[
G=\langle x_1,\ldots,x_k\mid r_1,\ldots,r_m\rangle
 =F(X)/\langle\!\langle r_1,\ldots,r_m\rangle\!\rangle,
\]
where the denominator is the normal closure of the relators. A group \(G\) is simple if \(G\neq\{1\}\) and its only normal subgroups are \(\{1\}\) and \(G\).

Order the formal alphabet \(X\cup X^{-1}\), and enumerate all finite words over this alphabet as \(w_0,w_1,\ldots\), first by length and then lexicographically, including the empty word. Write \(\overline{w_n}\) for the element represented by \(w_n\) in \(G\), and define the word-problem characteristic function by
\[
\chi_G(n)=
\begin{cases}
1,&\overline{w_n}=1_G,\\
0,&\overline{w_n}\neq 1_G.
\end{cases}
\]
Here a primitive recursive function means a function obtained from the zero, successor, and projection functions by composition and primitive recursion.

\textbf{Problem (F. B. Cannonito).} Construct an explicit finite presentation as above such that \(G\) is simple and \(\chi_G\) is not primitive recursive. Equivalently, require
\[
\forall f:\mathbb{N}\to\{0,1\}\text{ primitive recursive},\quad
\exists n\in\mathbb{N}\quad f(n)\neq\chi_G(n).
\]
The requested output is an explicit example, together with justification of these properties. The source does not prescribe a particular format for explicitness; specifying the finite presentation is the interpretation used here. The word encoding above makes the computational requirement precise.
\subsection{Short English statement}
Construct an explicit finitely presented, nontrivial simple group whose word problem cannot be decided by any primitive recursive function. The supplied record labels this problem open; that label is not an independently verified assessment.
\subsection{Sources}
\begin{itemize}
\item Local UnsolvedMath registry: supplied source record with permanent ID \(700419\), exact source problem code \texttt{catalogue entry}, titled \emph{Kourovka Notebook Problem 7.19}. Upstream revision: \texttt{d371ee917ac9647fbb06cd6fc529deec6f3dda82}.
\item \emph{The Kourovka Notebook (Unsolved Problems in Group Theory)}, Problem 7.19, originally in the 7th issue (1980); proposed by F. B. Cannonito. The supplied background identifies the 21st issue edition (2026), arXiv:1401.0300v48. The collection description names E. I. Khukhro and V. D. Mazurov as editors.
\item The supplied registry labels the problem \texttt{open} and records an update timestamp of \texttt{2026-10-16T14:00:00Z}. These are supplied registry metadata, not independent source verification or a new literature assessment. No separate literature-verification date is supplied.
\end{itemize}
\end{document}
